\subsection{Hyperparameter Tuning}

To optimize the performance of the \gls{xgBoost} classifier, hyperparameter tuning was performed using \texttt{\href{https://optuna.org/}{Optuna}} \citep{Akiba2019}. Hyperparameter optimization is necessary because the predictive performance of gradient-boosted decision trees depends strongly on the choice of training parameters, such as the learning rate, tree depth, regularization strength, and sampling ratios as described in \autoref{sec:xgBoost/regularization}. Rather than relying on manually selected values or an exhaustive grid search, \texttt{Optuna} was used to explore the search space efficiently and identify configurations that generalize well on the validation data.
\texttt{Optuna}'s optimization process is organized around the concept of a \emph{study}, which consists of multiple \emph{trials}. In each trial, \texttt{Optuna} samples one candidate set of hyperparameters, trains the model with these values, evaluates the resulting performance, and then uses the observed outcome to guide future trials.

A key advantage of \texttt{Optuna} is that it does not require all hyperparameter combinations to be defined and evaluated in advance. Instead, the search space is described programmatically inside an objective function. For each parameter, the trial object proposes values from predefined distributions, for example:
\begin{itemize}
    \item \texttt{suggest\_float(..., log=True)} for continuous parameters over several orders of magnitude,
    \item \texttt{suggest\_int(...)} for integer-valued parameters,
    \item \texttt{suggest\_categorical(...)} for discrete design choices.
\end{itemize}

Internally, \texttt{Optuna} uses adaptive search strategies to decide which parameter values should be explored next. In particular, it commonly uses the \emph{Tree-structured Parzen Estimator} (TPE), a Bayesian optimization method. TPE models the relationship between previously evaluated hyperparameter configurations and their objective values, and then proposes new configurations that are more likely to improve performance. Compared to random search or grid search, this makes the optimization substantially more sample-efficient, especially when the search space is large and contains a mixture of continuous, integer, and categorical variables.

Another important feature of \texttt{Optuna} is \emph{pruning}. Pruning allows unpromising trials to be stopped early if intermediate results indicate that they are unlikely to outperform the current best configuration. This reduces computational cost and allows the optimization to focus on more promising regions of the search space. In the present setup, the number of estimators was intentionally set to a large upper bound, while model selection was effectively controlled by early stopping during training. 
The study setup is  displayed in Listing~\ref{lst:optuna_study} in \autoref{app:methods}.
The evaluation metric \texttt{aucpr} was chosen because it is well suited for imbalanced binary classification tasks, as it emphasizes the precision--recall trade-off and is more informative than accuracy when the positive class is rare. The \texttt{hist} tree method was used to accelerate training by discretizing feature values into bins, and computation was performed on the CPU.

During optimization, each \texttt{Optuna} trial sampled one configuration from the above search space. The XGBoost model was then trained and evaluated on the validation data using the area under the precision--recall curve (\texttt{AUC-PR}) as the objective value. Based on the observed performance of completed trials, \texttt{Optuna} adaptively selected new hyperparameter configurations that were expected to yield better results.

This procedure has several advantages. First, it enables efficient exploration of a high-dimensional search space containing both continuous and categorical parameters. Second, it allocates more trials to promising regions of the search space rather than evaluating all configurations uniformly. Third, the use of log-scaled sampling for parameters such as \texttt{learning\_rate}, \texttt{min\_child\_weight}, \texttt{reg\_lambda}, and \texttt{reg\_alpha} is appropriate because the optimal values of these parameters often vary across several orders of magnitude.

In summary, \texttt{Optuna} was used as an adaptive hyperparameter optimization framework to tune the XGBoost classifier. By combining programmatically defined search spaces, Bayesian sampling through TPE, and efficient exploration of both continuous and discrete parameters, \texttt{Optuna} provided a flexible and computationally efficient way to identify a well-performing XGBoost configuration. The tuning focused on parameters controlling learning dynamics, tree complexity, feature and sample subsampling, regularization, and histogram discretization, while model quality was assessed with \texttt{AUC-PR} to reflect the requirements of the binary classification task.